% =====================================================================
%  Chapter 9 — The Verification, Part IV
% =====================================================================

\chapter{The Static Scan}
\label{ch:scan}

The scan is the cheapest thing one can do to a formalisation, and it is the
thing most often over-claimed. This chapter reports the scan results for the
submission, states precisely what they establish, and explains --- with a
worked example --- why they cannot settle the question.

\section{Method}The artifact was extracted locally and scanned in full. The scans were run
with \texttt{rg} (ripgrep\idx{static scan}) over the complete tree of 175{,}187 \texttt{.lean}
files, with \texttt{--no-ignore} so that nothing was skipped by any ignore
rule. No sampling was used. The patterns and results are in
\tabref{tab:scan}.

\begin{table}[H]
\centering
\caption{Static scan of the complete extraction.}
\label{tab:scan}
\begin{tabular}{@{}>{\RaggedRight\arraybackslash\hyphenpenalty=10000\exhyphenpenalty=10000}p{0.40\textwidth}r r@{}}
\toprule
\textbf{Pattern} & \textbf{This book's scan} & \textbf{Mirror's scan} \\
\midrule
\texttt{axiom} declarations        & 0 & 0 \\
\texttt{sorry}                     & 0 & 0 \\
\texttt{admit}\idx{admit}                     & --- & 0 \\
\texttt{native\_decide}            & --- & 0 \\
\texttt{unsafe} / \texttt{implemented\_by} / \texttt{@[extern]} & --- & 0 \\
\texttt{opaque} declarations       & 0 & 0 \\
\code{set\_option maxHeartbeats 0}\idx{maxHeartbeats} & --- & 4 files \\
custom \texttt{macro} / \texttt{elab} / \texttt{syntax} & --- & 22 files \\
\texttt{addDecl}\idx{addDecl}                   & --- & 1 file \\
\bottomrule
\end{tabular}
\end{table}

The two scans agree on everything they both cover. This is a good sign: two
independent scans, by different parties with different tooling, produce the
same numbers on the two checks that matter most.

\section{What the scan establishes}Three things, and they are real.

\begin{quote}
\textbf{There is no unproved axiom.} The development does not begin by
declaring the Riemann hypothesis to be true. Whatever else is going on, it is
not the trivial kind of fraud.

\textbf{There is no admitted hole.} No \texttt{sorry}, no \texttt{admit}, no
\texttt{native\_decide} whose soundness depends on the compiler. Every theorem
in the development is, in Lean's sense, closed.

\textbf{There is no custom kernel.} No \texttt{unsafe} declarations, no
\texttt{implemented\_by}, no \texttt{opaque} declarations, no
\texttt{addDecl} outside the single file the mirror flags. The kernel is
checking what it thinks it is checking.
\end{quote}

\section{What the scan does not establish}And here is the part that matters. A scan is a search for a \emph{marker}. The
markers above are markers of a specific, common failure mode. The failure mode
they detect is not the one most likely to be present here.

\begin{deeper}{A worked example of the failure mode the scan cannot see}
Suppose a formalisation wants to prove that the Riemann zeta function has no
zeros in $\mathrm{Re}\,  s > 1/2$. There is a legitimate-looking way to organise this:

\begin{leanquote}
-- A hypothesis bundle, named for a classical theorem
structure WeilBoundHolds (f : Complex → Complex) : Prop where
  positivity : ∀ (g : Complex → Complex), admissible g →
    sumZeros f g >= 0

-- The "theorem" that does the work
theorem zetaHasNoZerosRightOfHalf (h : WeilBoundHolds riemannZeta) :
    ∀ s : Complex, 1/2 < s.re → riemannZeta s ≠ 0 := by
  intro s hs
  exact h.positivity _ admissible s hs
\end{leanquote}

This compiles. The axiom list is clean. There is no \texttt{sorry}. The
structure is named after a classical theorem. And the development now has a
theorem called \code{zetaHasNoZerosRightOfHalf} whose proof is one line long.

But what has been shown is: \emph{if the Weil positivity statement holds for
$\zeta$, then $\zeta$ has no zeros to the right of the critical line}. The
content is entirely in the hypothesis, and the hypothesis has been given a
friendly name.

To detect this, the reader must do the thing the scan cannot do: \emph{read
the statement of the hypothesis and ask whether it is the classical theorem.}
In the case above, the statement of the hypothesis is \emph{not} the classical
Weil criterion --- the classical criterion is a positivity condition on a
quadratic form $\sum_\rho \hat g(\rho)\hat g(1-\rho) \ge 0$, and the hypothesis
above says something involving a sum over zeros of an arbitrary function, which
is a different statement, stronger in some directions and meaningless in
others.
\end{deeper}

This is the failure mode that matters for the artifact under examination, and
it is why Chapter~8 spends its effort on \emph{statements} rather than on
proof terms.

\begin{keytakeaways}

\item Two independent scans agree: zero axioms, zero \texttt{sorry}, zero
\texttt{native\_decide}, zero custom kernel escapes.
\item A scan searches for markers of a specific failure mode. The failure mode
most likely to be present in a large, sophisticated development is a
\emph{weakened hypothesis statement wearing a classical name}, and that
failure mode leaves no marker at all.
\item A theorem whose proof is one line long, whose hypothesis is a structure
with a classical name, and whose axiom list is clean, has told you nothing
until you have read the structure's fields against the published theorem.
\end{keytakeaways}

\section{The scan is nevertheless worth doing, and here is why}It would be wrong to conclude that the scan was wasted. Three reasons.

\begin{enumerate}
\item \textbf{It is a precondition, and it was not a foregone conclusion.} It
is entirely possible --- and it has happened in published work --- for a
formalisation to smuggle in an axiom, or to use \texttt{native\_decide}, or
to depend on a custom kernel extension. The scan rules all of these out at
zero cost.

\item \textbf{It bounds the search space for the human work.} A reader who
knows the axiom list is clean and there is no \texttt{sorry} knows that the
remaining work is entirely in the statements, and can direct their attention
there. This is not a small saving on a 175-thousand-file development.

\item \textbf{It is reproducible, and reproducibility is the point.} Any
reader with the archive and twenty minutes can run these three commands and
get the same answers. A book that reports its scan results, with the commands
and the tree size, is offering the reader a check rather than a conclusion.
\end{enumerate}

\section{One specific thing to look at first}The mirror's scan flags three items that this book does not treat as evidence
of anything, but which are the right place to start a human audit, for one
reason: they are the places where the \emph{elaborator} rather than the
kernel does something unusual.

\begin{description}[style=nextline,leftmargin=0pt,labelwidth=0pt,labelsep=0pt,font=\sffamily\bfseries\small,
  itemsep=5pt]
\item[\texttt{addDecl}, one file.] \texttt{addDecl} injects a declaration into
the environment directly, bypassing the elaborator's usual checks. It is used
legitimately --- for instance to register a \texttt{notation} or an
\texttt{instance} --- but a reader who finds it should read the surrounding
code and satisfy themselves that it is not defining something that should have
been a theorem.
\item[Custom \texttt{macro}/\texttt{elab}/\texttt{syntax}, 22 files.] A
custom elaborator macro can expand a proof obligation into something trivial.
The defence against this is the same as the defence against any other form of
the attack: read it. But it must be read.
\item[\code{set\_option maxHeartbeats 0}, four files.] This disables the
elaborator's heartbeat limit for a declaration, which usually means the
declaration is expensive to elaborate --- a large \texttt{simp} set, a big
\texttt{decide}, or an exhaustive search. It is not a soundness problem. It is
a place where the proof term might be large and might not have been inspected
closely.
\end{description}

None of these is a smoking gun. All three are reasonable things for a
175-thousand-file development to contain. They are listed here because the
point of listing them is that they are checkable, and because a book that
reported only the reassuring results would not be doing its job.

\section{The scan this book could not run}There is one check that would strengthen everything above, and this book
could not perform it: the axiom check itself.

Running \code{\#print axioms riemannHypothesis} requires a working Lean
installation with the project compiled, because the command has to elaborate
the theorem's transitive dependencies. The static scan establishes that
there is no \texttt{axiom} keyword anywhere in the source, which is
\emph{almost} the same thing --- but not quite. A dependency could in principle
introduce an axiom outside the source tree, in Mathlib itself, and the
development pins Mathlib at a specific revision precisely so that this can be
checked.

In practice the risk is negligible: Mathlib's axiom set is the three standard
ones, and this is verified as a matter of course in Mathlib's own CI. But
``verified as a matter of course'' is not the same as ``verified for this
theorem in this development'', and the difference is exactly the kind of
difference this book has been insisting on since Chapter~6.

Chapter~10 reports the outcome of the build attempt made for this book, and
Chapter~14 gives the command.
